\documentclass[10pt,a4paper]{article}
\usepackage[utf8]{inputenc}
\usepackage[a4paper,left=3cm,right=2cm]{geometry}
\usepackage{amsmath}
\usepackage[thmmarks, amsmath, thref]{ntheorem}
\usepackage{amsfonts}
\usepackage{amssymb}
\usepackage{fancyhdr}
\pagestyle{fancy}
\usepackage{anysize}
\usepackage{graphicx}
\usepackage{float}
\usepackage{hyperref}
\usepackage{enumitem}
\newtheorem{theorem}{Theorem}
\theoremsymbol{\ensuremath{\blacksquare}}
\newtheorem*{solution}{Solución:}
\usepackage{float}
\usepackage{xcolor}

% Margenes y formalidades

\textwidth 6.25in % Margen Derecho y ancho del texto
\textheight 8.5in
\oddsidemargin 0in % Margen Izquierdo
\headheight 0in % Largo del texto
\leftmargin 1in
\parindent 0em % Salto entre la indentación (Sangría)https://www.overleaf.com/4519127669sdpptkhvsbfn
\topmargin -0.5in % Margen Superior
\parskip 0.5ex % Salto entre parrafos
\baselineskip 1pt



\lhead{
\includegraphics[scale=0.09]{Figuras/logoUSM.png}} % texto izquierda de la cabecera
%\chead{}                                      % texto centro de la cabecera
\rhead{Universidad Técnica Federico Santa María\\Ingeniería Comercial } % número de página a la derecha
\renewcommand{\headrulewidth}{0.4pt} % grosor de la línea de la cabecera
\renewcommand{\footrulewidth}{0.4pt} % grosor de la línea del pie

\begin{document}
\vspace{5cm}



\begin{center}
\begin{large}
\textbf{\\CERTAMEN 3}
\end{large}\\
\end{center}

\begin{center}
\begin{large}
\textbf{Econometría ICS-294}
\end{large}\\
\vspace{0.2cm}
II trimestre 2025\\
Profesor: Pedro Comte Hidalgo
\end{center}

\textbf{Instrucciones:} 

\begin{itemize}

    \item Usted dispone de 100 minutos  para intentar los 100 puntos del certamen.
    \item Cuide la presentación y redacción de sus preguntas.
    \item No puede utilizar celular, computador o tablet, pero sí una calculadora TI.
     
\end{itemize}

\section{Preguntas Cortas (40 Puntos)}

\begin{enumerate}

\item (8 Puntos) ¿Que es la heterocedasticidad y por qué en presencia de ella la inferencia estadística deja de ser consistente?
\\
\textbf{Solución}
\\
La heterocedasticidad ocurre cuando la varianza de los errores no es constante a lo largo de las observaciones. La inferencia estadística deja de ser válida porque los errores estándar calculados de manera tradicional resultan incorrectos, de modo que los estadísticos t y F pierden validez.

\item (8 Puntos) Suponga un modelo que mide la probabilidad de cometer un asesinato dado que la persona realizó el servicio militar. Sabiendo que existe endogeneidad, ¿cuál podría ser una variable instrumental válida para este problema y por qué? \\
\textbf{Solución} 
\\

Una variable instrumental en este problema es si la persona fue escogida aleatoriamente para realizar el servicio militar y no lo hizo de manera voluntaria. En este caso se cumplen las dos condiciones que debe tener una variable instrumental: 

\begin{enumerate}
    \item \textbf{Relevancia}: la covarianza entre la variable instrumental $Z$ y la variable explicativa $X$ es distinta de cero; en este contexto, haber sido sorteado está correlacionado con haber realizado el servicio militar. 
    \item \textbf{Exogeneidad}: el sorteo es aleatorio y, por lo tanto, no tiene correlación con el término de error ni con otras características individuales relacionadas con la probabilidad de cometer un asesinato.
\end{enumerate}


\item (8 Puntos) ¿Cuál es la diferencia entre multicolinealidad perfecta y alta multicolinealidad, y como afectan al calculo de los regresores mediante MCO?\\
\textbf{Solución}
\\
La \textit{multicolinealidad perfecta} se presenta cuando una variable explicativa es una combinación lineal exacta de otras variables del modelo. En este caso, la matriz $X'X$ es singular y no puede invertirse, por lo que los estimadores de MCO no pueden calcularse: el modelo no es identificable.  

En cambio, la \textit{alta multicolinealidad} ocurre cuando las variables explicativas están fuertemente correlacionadas, pero no de manera exacta. Aquí los estimadores de MCO sí pueden calcularse, pero presentan una varianza elevada, lo que se traduce en errores estándar grandes y dificulta la inferencia estadística, ya que los coeficientes pueden aparecer como no significativos aun cuando la variable sea relevante, y dificultando su interpretación.

\item (8 puntos) Explique brevemente por qué una variable que sigue una caminata aleatoria (\textit{random walk}) no es estacionaria.\\
\textbf{Solución}\\
Una caminata aleatoria no es estacionaria porque su media depende de la condición inicial y su varianza crece sin límite a medida que aumenta el tiempo. Esto implica que la distribución de la serie cambia en el tiempo, por lo que no cumple con los requisitos de estacionariedad.

\item (8 puntos) ¿Qué efecto tiene aplicar primeras diferencias a una serie con raíz unitaria?
\\
\textbf{Solución}
\\
Aplicar primeras diferencias a una serie con raíz unitaria (por ejemplo, un random walk) elimina la no estacionariedad y transforma la serie en estacionaria.

\end{enumerate}

\section{Ejercicios (60 Puntos)}

\begin{enumerate}

\item (20 Puntos) Se dispone de una muestra de 500 observaciones con la cual se estimó el siguiente modelo de regresión lineal múltiple:
$$Salario_i = \beta_0 + \beta_1\,\textit{Educ}_i + \beta_2\,\textit{Exper}_i + u_i$$
Se desea contrastar el supuesto de heterocedasticidad y, para ello, se aplicará el método de White (no simplificado).
\begin{enumerate}
    \item (6 Puntos) Plantee la regresión auxiliar para el método de White (no simplificado).\\
    \textbf{Solución}\\
    $$\hat{u}_i^{\,2} = \alpha_0 + \alpha_1 {Educ}_i + \alpha_2 {Exper}_i + \alpha_3 Educ_i\cdot Exper_i + \alpha_4 Educ_i^2+\alpha_5 Exper_i^2+v_i$$
    
    \item (4 puntos) Formule la hipótesis nula y la hipótesis alternativa para contrastar la presencia de heterocedasticidad.\\
    \textbf{Solución}\\
    $$H_0:\ \alpha_1 = \alpha_2 =\alpha_3 =\alpha_4 = \alpha_5=0$$
    $$H_1:\ \alpha_1 \neq 0 \ \text{o}\ \alpha_2 \neq 0 \ \text{o} \ \alpha_3 \neq 0 \ \text{o} \ \alpha_4 \neq 0 \ \text{o} \ \alpha_5 \neq 0$$
    
    \item (10 Puntos) Si el $R^2$ de la regresión auxiliar es $0.03$, contraste, al 1\% de significación, la existencia de heterocedasticidad.\\
    \textbf{Solución}\\
    $$LM \;=\; n R^2_{\text{aux}} \;=\; 500 \times 0.03\;=\; 15 \;\sim\; \chi^2_{(5)}$$
    El valor crítico de $\chi^2_{(2)}$ al 1\% es $15.086$. Como
    $$LM \;>\; \chi^2_{0.95,\,2} \;\; \Longrightarrow \;\; 15 \;<\; 15.086,$$
    No se rechaza la hipótesis nula y se concluye que existe homocedasticidad.
\end{enumerate}

\item (20 Puntos) Un investigador desea analizar los determinantes del rendimiento académico universitario mediante la siguiente regresión múltiple:
\[
\text{PROM} = \beta_0 + \beta_1 \,\text{HES} + \beta_2 \,\text{ASIST} + \beta_3 \,\text{TE} + u,
\]
donde:
\begin{itemize}
  \item \textbf{PROM}: promedio final del estudiante,
  \item \textbf{HES}: horas de estudio semanal,
  \item \textbf{ASIST}: asistencia a clases (en \%),
  \item \textbf{TE}: número de trabajos entregados.
\end{itemize}

Se dispone de la siguiente \textbf{matriz de correlación} entre las variables explicativas \(\{\text{HES}, \text{ASIST}, \text{TE}\}\), respectivamente:
\[
\begin{bmatrix}
1 & 0.90 & 0.85 \\
0.90 & 1 & 0.70 \\
0.85 & 0.70 & 1
\end{bmatrix}
\]

Para evaluar la posible presencia de multicolinealidad, se estimaron las regresiones auxiliares de cada regresor sobre los otros dos, obteniéndose los siguientes coeficientes de determinación:
\[
R^2_{\text{HES}\mid \text{ASIST},\text{TE}} = 0.91,\quad
R^2_{\text{ASIST}\mid \text{HES},\text{TE}} = 0.75,\quad
R^2_{\text{TE}\mid \text{HES},\text{ASIST}} = 0.80.
\]
\begin{enumerate}
  \item (6 puntos) Examine la matriz de correlación e indique si existen indicios de multicolinealidad entre las variables explicativas. Señale qué pares presentan correlaciones particularmente altas y por qué esto puede ser problemático.\\
  \textbf{Solución}
  \\
  Las variables HES con ASIST tienen la mayor correlación (0.9), pero tambien HES con TE tienen alta correlación (0.85), lo que puede complicar las interpretaciones de los regresores y tambien aumenta su varianza.
  
  \item (14 puntos) Calcule el \textit{Factor de Inflación de la Varianza} (VIF) e interprete los VIF obtenidos. ¿Existe evidencia de multicolinealidad severa y sobre que variable? Proponga al menos una estrategia para mitigar el problema.\\
  \textbf{Solución}
  \\
    $$VIF_{HES}=\frac{1}{1-0.91}=11.1$$
    $$VIF_{ASIST}=\frac{1}{1-0.75}=4$$
    $$VIF_{TE}=\frac{1}{1-0.8}=5$$
    Si hay evidencia de multicolinealidad sobre la variable HES. Se propone eliminar del modelo la variable HES para poder reducir la varianza de los estimadores y facilitar la interpretación. Otra alternativa es aplicar una transformación lineal, como aplicar logaritmo natural, pero en este caso se tendría que volver a medir su VIF.
 \end{enumerate}

\item (20 Puntos) Se tiene la siguiente grafica del PIB de chile en Millones de Euros, desde 1999 a 2024:
\begin{center}
    \begin{figure}[H]
        \centering
        \includegraphics[width=0.6\linewidth]{certamenes/Certamenes T2/imagen pib.png}
        \caption{PIB anual Chile}
        \label{fig:enter-label}
    \end{figure}
\end{center}
\begin{enumerate}
    \item (6 Puntos) Que puede opinar de la grafica observada. ¿Existe tendencia? ¿Que modelo de regresión lineal en series de tiempo plantearía usted para el PIB anual de Chile? 
    \\
    \textbf{Solución}
    \\
Se puede observar que la grafica del pib tiene una tendencia al crecimiento con pequeñas fluctuaciones a la baja aprox. cada 8 años. Como el PIB tiene tendencia, y sus valores siguen creciendo, se deberia medir en diferencias, por lo que el modelo que yo plantearía sería:
$$\Delta PIB_t = \beta_0 + \phi PIB_{t-1} + \alpha t + u_t$$
Siendo $\phi=\rho-1$
    \item (4 Puntos) Plantee la hipotesis nula y alternativa del modelo antes planteado para testear si existe raiz unitaria. 
    \\     \textbf{Solución}
    $$H_0 : \phi = 0$$
    $$H_1 : \phi < 0$$
    \\
    \item (10 puntos) Según los siguientes resultados del modelo de regresión lineal en series de tiempo, ¿existe raíz unitaria al 5\% de significancia? Compruebe usando test de Dickey-Fuller.
    \begin{center}
    \begin{figure}[H]
        \centering
        \includegraphics[width=0.6\linewidth]{certamenes/Certamenes T2/rlst.png}
        \caption{Resultados Regresión}
        \label{fig:enter-label}
    \end{figure}
\end{center}
\end{enumerate}
\\ \textbf{Solución}
\\
Estadistico DF:
$$DF_{est} = \frac{-0.5208}{0.1777}=-2.931$$
Según la tabla el estadistico critico es:
$$DF_{Critico}=-3.6$$
Como
$$|-2.931|<|-3.6|$$ entonces no se rechaza $H_0$ y hay evidencia de que existe raiz unitaria en el proceso del PIB.

\end{enumerate}


\section*{Formulario}
\begin{itemize}

\item Tabla $\chi^2$
\begin{center}
    \begin{tabular}{|c|c|c|c|}
    \hline
        df & 0.99 & 0.95 & 0.90 \\ 
        \hline
        1 &  6.635 & 3.841 & 2.706\\
        \hline
        2 &  9.210 & 5.991 & 4.605 \\
        \hline
        3 & 11.345 & 7.815 & 6.251  \\
        \hline
        4 &  13.277 & 9.488 & 7.779 \\
        \hline
        5 &  15.086 & 11.070 & 9.236 \\
        \hline
        6 &  16.812 & 12.592 & 10645 \\
        \hline
    \end{tabular}
\end{center}

\item Multiplicado de Lagrange

$$\color{red}{LM = nR^2 \sim \chi^2}$$

\item Test F para regresión auxiliar.
\begin{center}
 {  \color{red}     \(F=\dfrac{R_{\text{aux}}^{2}/q}{(1-R_{\text{aux}}^{2})/(n-k-1)}
          \sim F_{q,\;n-k-1}\).}
\end{center}



\item Tabla de test D-F
\begin{center}
    \begin{figure}[H]
        \centering
        \includegraphics[width=0.7\linewidth]{sem14/tabla.png}
    \end{figure}
\end{center}

\end{itemize}


\end{document}